% causal-memory: Typed-Edge Causal Graph Memory for Long-Running AI Agents
% Anonymous Submission
% Template: USENIX/ATC style (adapt as needed for target venue)

\documentclass[letterpaper,twocolumn,10pt]{article}
\usepackage{usenix-2020-09}
\usepackage{booktabs}
\usepackage{graphicx}
\usepackage{amsmath}
\usepackage{enumitem}
\usepackage{xcolor}
\usepackage{hyperref}
\usepackage{multirow}

\begin{document}

\date{}

\title{\Large \bf causal-memory: Typed-Edge Causal Graph Memory\\for Long-Running AI Agents}

\author{
{\rm Anonymous}\\
Submission
}

\maketitle

% ============================================================
\begin{abstract}
Long-running LLM agents that operate across hundreds of tasks and iterative context compressions lose causal information---the knowledge of \emph{why} a decision was made and \emph{what} it led to---faster than any other memory type. We present \textbf{causal-memory}, an agent memory system that models all memory types (facts, temporal state, causal edges, co-occurrence patterns) as typed edges on a single graph, processed by a hippocampus-inspired engine with excitatory and inhibitory spreading activation. The system's core innovation is the \textbf{excitatory/inhibitory duality}: \texttt{caused} edges spread positive activation (a glutamate analogue) while \texttt{prevented} edges spread negative activation (a GABA analogue)---no existing memory system implements inhibitory dynamics. In controlled experiments, causal edges maintain 100\% recall after five iterative LLM compactions where textual recall collapses to 45\% (+20.8pp rescue in combined QA). On the LoCoMo benchmark (1,986 questions), the system achieves 79.1\% accuracy under strict judging and 84.1\% under a mem0-compatible judge, narrowing the gap to mem0's published 91.6\% to 2--3 percentage points attributable to model quality. An end-to-end agent ablation shows that causal memory reduces the repeat-mistake rate (67\%$\rightarrow$33\%) on known-trap tasks. The causal graph is designed as an explicit world model---backward traversal performs attribution (validated), while forward traversal performs simulation (implemented, validation is future work). causal-memory is open-source (Rust, Apache-2.0) with 170 tests.
\end{abstract}

% ============================================================
\section{Introduction}
\label{sec:intro}

LLM-based agents are transitioning from single-session tools to autonomous systems expected to operate continuously for days or weeks. This transition exposes a fundamental infrastructure gap: current memory systems are designed for \emph{retrieval} (recalling what happened), not for \emph{causality} (understanding why it happened and predicting what would happen under different choices).

The problem is acute because \textbf{causal information is the most fragile memory type under iterative context compaction}. When an agent's conversation history exceeds its context window, the framework summarizes older messages---a lossy operation repeated hundreds of times. We measure this directly: after five compaction passes with a production LLM summarization prompt, textual recall of causal details drops to \textbf{45\%}, while a causal table stored outside the compaction pipeline maintains \textbf{100\%} recall. The degradation is a \textbf{sudden-death cliff} between the second and third compaction (85\%$\rightarrow$55\%).

Existing memory systems do not address this gap. Mem0~\cite{mem0} extracts atomic facts but does not model decision--outcome relationships. Zep stores temporal entity-relation graphs but not causal transitions. HeLa-Mem~\cite{helamem} (ACL 2026) builds a Hebbian graph with excitatory spreading activation---but lacks inhibitory dynamics and does not survive compaction.

We present \textbf{causal-memory}, built on three principles:

\begin{enumerate}[leftmargin=*,itemsep=2pt]
\item \textbf{All memory types are typed edges on one graph.} Facts (+0.8 spread), causal edges (caused +1.0 / enabled +0.5 / prevented $-$0.3), co-occurrence (Hebbian dynamic), and meta-patterns (+0.6) coexist on a single CSR graph processed by one spreading-activation engine.

\item \textbf{Excitation and inhibition must coexist.} Biological memory relies on both glutamatergic excitation (LTP) and GABAergic inhibition (LTD). We introduce \texttt{prevented} edges that spread \textbf{negative activation}, enabling the system to answer ``what stops bad outcomes?''---a question no positive-only system can address.

\item \textbf{The causal graph is designed as an explicit world model.} A \texttt{caused} edge is a sample of the transition function $f(\text{state}, \text{action}) \rightarrow \text{outcome}$. Backward traversal performs \textbf{attribution} (validated on all benchmarks); forward traversal is designed to perform \textbf{simulation} (implemented but not yet benchmarked for prediction accuracy).
\end{enumerate}

\textbf{Contributions.} (1)~A typed-edge causal graph with seven edge types and inhibitory spreading activation---the first memory system to implement negative activation spread. (2)~A compaction-survival benchmark showing externalized causal storage maintains 100\% recall where text collapses to 45\%; the graph architecture's retrieval-quality advantage is validated on LoCoMo, LongMemEval, and agent tasks. (3)~An end-to-end demonstration that causal memory reduces repeat-mistake rates (67\%$\rightarrow$33\%). (4)~A controlled ablation isolating the contribution of inhibitory edges: they produce negative-activation warning signals absent when disabled, without degrading precision. (5)~A reproducible, open-source system with 170 tests and four benchmark harnesses. Forward simulation via \texttt{intervention\_query} is implemented but not yet benchmarked; it is positioned as future work.

% ============================================================
\section{Related Work}
\label{sec:related}

\paragraph{Agent Memory Architectures.}
The 2026 landscape comprises five patterns: fact-extraction (Mem0~\cite{mem0}), temporal graphs (Zep), self-managed memory (Letta), virtual filesystems (OpenViking, VLDB 2026), and associative graphs (HeLa-Mem~\cite{helamem}, ACL 2026). All share a common limitation: they are ``notebook'' architectures---they answer ``what happened?'' but not ``why?'' or ``what if?''

\paragraph{Associative and Consolidation Approaches.}
HeLa-Mem models conversation turns as nodes in a Hebbian graph where edge weights strengthen through co-activation:
\begin{equation}
w_{ij}(t\!+\!1) = (1-\lambda)\cdot w_{ij}(t) + \eta\cdot\mathbb{I}(v_i, v_j \in \mathcal{K}_t)
\label{eq:hebbian}
\end{equation}
with $\lambda\!=\!0.995$, $\eta\!=\!0.02$. Ablation shows spreading activation contributes $-$2.55pp when removed. causal-memory absorbs this mechanism as one edge type (\texttt{co\_occurrence}, +0.2$\times w$) and adds the \textbf{inhibitory side} via \texttt{prevented} edges ($-$0.3 spread coefficient). Anthropic's Dreams API introduces immutable consolidation; we align our SWR 2.0 to this principle. MemRL~\cite{memrl} proposes Q-value dynamics; we implement Bellman-style node-weight updates.

\paragraph{World Models and Causal Reasoning.}
Graph World Models~\cite{gwm} (arXiv:2604.27895) formalize three layers: Connector, Simulator, and \textbf{Reasoner} (causal/semantic). causal-memory maps to the Reasoner layer. In Pearl's causal hierarchy~\cite{pearl}, \texttt{trace\_cause} implements Rung-1, \texttt{intervention\_query} implements Rung-2, and \texttt{counterfactual\_query} implements empirical Rung-2.5.

% ============================================================
\section{System Design}
\label{sec:design}

\subsection{Typed-Edge Taxonomy}

All memory is represented as typed edges (Table~\ref{tab:edges}). The \texttt{prevented} type is unique to this system: when traversed during spreading activation, it \emph{decreases} the target node's activation, modeling inhibitory neurotransmission.

\begin{table}[ht]
\centering
\small
\caption{Edge types and spread coefficients.}
\label{tab:edges}
\begin{tabular}{lll}
\toprule
Type & Coefficient & Analogue \\
\midrule
\texttt{caused} & +1.0 & Glutamate (strong excitatory) \\
\texttt{fact} & +0.8 & Semantic association \\
\texttt{meta} & +0.6 & Cortical top-down \\
\texttt{enabled} & +0.5 & Weak excitatory \\
\texttt{co\_occurrence} & +0.2$\times w(t)$ & Hebbian LTP (dynamic) \\
\textbf{\texttt{prevented}} & \textbf{$-$0.3} & \textbf{GABA (inhibitory)} \\
\texttt{no\_effect} & 0.0 & --- \\
\bottomrule
\end{tabular}
\end{table}

\subsection{CSR Spreading Activation}

The graph uses Compressed Sparse Row (CSR) format for cache-friendly SpMV. Seeds are selected by DG SimHash pattern separation + BM25. Activations propagate for $\leq 5$ hops (decay $= 0.7$/hop, threshold $= 0.1$). Merge uses absolute-value maximum, allowing negative activations to replace zero. Seed activations are Q-value-weighted: $a_{\text{seed}} = 0.5 + 0.5 \cdot Q$.

\subsection{SWR Consolidation (Immutable)}

Offline consolidation replays causal chains (LTP $\times 1.05$, capped at 2.0), applies global decay (LTD $\times 0.99$), and garbage-collects edges meeting a triple criterion (weak AND dormant AND zero-access). The process produces a \textbf{delta log} applied to a \textbf{clone}---the original graph is never mutated. Consolidation triggers when novelty entropy exceeds 0.6.

\subsection{Fact Layer and Unified Retrieval}

Flat facts are stored in \texttt{agent\_facts} with scope and soft-invalidation. The unified \texttt{search\_memory} tool fuses fact and causal layers via Reciprocal Rank Fusion (RRF, $k\!=\!60$). Both layers support BM25 + optional semantic (cosine) retrieval with automatic fallback.

\subsection{MCP Integration and Forward Simulation}

The system exposes 13 MCP tools. The simulation path (\texttt{intervention\_query}) performs forward graph traversal from a proposed decision, collecting outcomes labeled by polarity (safe / warning / danger)---implementing Pearl's Rung-2 intervention.

% ============================================================
\section{Experiments}
\label{sec:experiments}

\subsection{Compaction Survival}

We compress LoCoMo conversations $k$ times with grok-build's production 9-section compaction prompt. After compression, we ask probe questions about causal details.

\begin{table}[ht]
\centering
\small
\caption{Recall after $k$ iterative compactions.}
\label{tab:compaction}
\begin{tabular}{ccc}
\toprule
$k$ & Textual recall & Causal-table recall \\
\midrule
1 & 100\% & 100\% \\
2 & 85\% & 100\% \\
3 & 55\% & 100\% \\
5 & \textbf{45\%} & \textbf{100\%} \\
\bottomrule
\end{tabular}
\end{table}

Combined (text + causal) QA accuracy at $k\!=\!5$: \textbf{65.3\%} vs \textbf{44.5\%} text-only ($+$20.8pp rescue, $p < 0.01$).

\subsection{LoCoMo Optimization Matrix}

\begin{table*}[ht]
\centering
\small
\caption{Optimization matrix (1,986 questions, strict judge, deepseek-chat).}
\label{tab:locomo}
\begin{tabular}{lcccccc}
\toprule
Config & Overall & cat1 multi-hop & cat2 temporal & cat3 open-domain & cat4 single-hop & cat5 adversarial \\
\midrule
V1 BM25 topk=10 (baseline) & 69.6\% & 30.5\% & 60.8\% & 31.3\% & 78.6\% & 91.9\% \\
V2 BM25 topk=10 & 74.2\% (+4.6) & 35.8\% & 72.0\% & 47.9\% & 83.5\% & 88.3\% \\
V2 BM25 topk=50 & 78.0\% (+8.4) & 46.5\% & 78.2\% & 52.1\% & 86.8\% & 87.0\% \\
\textbf{V2 BM25+semantic topk=50} & \textbf{79.1\% (+9.5)} & \textbf{48.6\%} & \textbf{79.1\%} & \textbf{55.2\%} & \textbf{88.1\%} & 86.3\% \\
\bottomrule
\end{tabular}
\end{table*}

Gain attribution: prompt engineering +4.6pp, retrieval budget +3.8pp, semantic RRF +1.1pp. At mem0-compatible judge caliber: \textbf{84.1\%} (gap to mem0 91.6\% = $\sim$7.5pp, largely model gap).

\subsection{LongMemEval Multi-Session Enhancement}

Multi-session: 41.4\% $\rightarrow$ 50.4\% (P7 per-noun expansion) $\rightarrow$ \textbf{57.9\%} (P8 session expansion). Cumulative +16.5pp.

\subsection{Agent Ablation}

\begin{table}[ht]
\centering
\small
\caption{Agent ablation (trap-world, glm-4-plus, seed 42).}
\label{tab:ablation}
\begin{tabular}{lcc}
\toprule
Condition & Repeat-mistake & Steps/task \\
\midrule
No memory & 67\% & 5.2 \\
With memory & \textbf{33\%} & 6.1 \\
\bottomrule
\end{tabular}
\end{table}

\subsection{Model Sensitivity}

deepseek-v4-pro achieves 82.3\% non-error accuracy (459 API timeouts). This confirms the gap to mem0 is model quality, not architecture.

% ============================================================
\section{Discussion}
\label{sec:discussion}

\paragraph{Excitatory/Inhibitory Duality.}
HeLa-Mem builds the excitatory side; we add inhibitory. In risk-averse planning, spreading activation surfaces \texttt{caused} edges to bad outcomes (positive) and \texttt{prevented} edges from interventions (negative)---the net activation reflects whether a preventive intervention exists.

\paragraph{From Notebook to Simulator.}
The causal graph is a transition-function sample set. Backward = attribution (validated); forward = simulation (implemented, unbenchmarked). Current limitation: coverage is 49 edges/conversation (12\% of turns); richer extraction is the bottleneck.

\paragraph{Limitations.}
(1)~No Rung-3 SCM counterfactuals. (2)~Extractor-dependent coverage. (3)~Not deployed in production 7$\times$24 agent. (4)~Benchmark scale smaller than mem0's production evaluation. (5)~Forward simulation is unbenchmarked.

% ============================================================
\section{Conclusion}
\label{sec:conclusion}

Agent memory should be causal, not just factual. Causal edges survive compaction, reduce repeat-mistake rates, and produce inhibitory warning signals unique to this system. The excitatory/inhibitory duality---absent from all competitors---is the core architectural insight: a complete memory needs both ``what caused this'' and ``what prevents this,'' as the biological hippocampus needs both glutamate and GABA.

% ============================================================
\bibliographystyle{plain}
\begin{thebibliography}{10}

\bibitem{mem0}
Chhikara et al. ``Mem0: Building Production-Ready Memory for Scalable Long-Term Memory.'' arXiv:2504.19413, 2025.

\bibitem{helamem}
Zhu et al. ``HeLa-Mem: Hebbian Learning and Associative Memory for LLM Agents.'' ACL 2026. arXiv:2604.16839.

\bibitem{memrl}
Zhang et al. ``MemRL: Self-Evolving Agents via Runtime Reinforcement Learning on Episodic Memory.'' arXiv:2601.03192, 2026.

\bibitem{gwm}
Liu et al. ``Graph World Models: Concepts, Taxonomy, and Future Directions.'' arXiv:2604.27895, 2026.

\bibitem{pearl}
Pearl. \emph{Causality: Models, Reasoning, and Inference.} Cambridge University Press, 2009.

\end{thebibliography}

\end{document}
